%% kapa: exam=1 round=31 year=2020 period=2 form=A subject=appraisal_law subject_ko=감정평가관계법규 seq=15 number=15
%% source: Q-Net 감정평가사 기출문제 (cst00309) artlSeq=5207219 file=extracted/1008/5207219/A형/A형/2020년도 제31회 감정평가사 1차 시험_2교시_A형.pdf
%% conversion: kordoc 4.24.51 -> tools/gen1.py (auto); review=proofread-sonnet; flags=none
\begin{question}{15}
감정평가 및 감정평가사에 관한 법령상 감정평가사의 권리와 의무 등에 관한 설명으로 옳지 않은 것은?
\begin{choices}
  \item 감정평가사합동사무소에 두는 감정평가사의 수는 2명 이상으로 한다.
  \item 감정평가사사무소의 폐업신고와 관련하여 국토교통부장관은 폐업신고의 접수업무를 한국감정평가사협회에 위탁한다.
  \item 감정평가업자는 소속 감정평가사의 고용관계가 종료된 때에는 한국감정원에 신고하여야 한다.
  \item 감정평가업자는 고의 또는 중대한 과실로 잘못된 평가를 하여서는 아니 된다.
  \item 감정평가업자가 감정평가를 하면서 고의 또는 과실로 감정평가 당시의 적정가격과 현저한 차이가 있게 감정평가를 함으로써 선의의 제3자에게 손해를 발생하게 하였을 때에는 그 손해를 배상할 책임이 있다.
\end{choices}
\end{question}
\answer{3}
